\section{Trace Schema and Validation Model}

ROK treats traceability as a first-class design requirement. All protocol interactions emit structured, append-only trace events that conform to an explicit, versioned schema.

\subsection{Append-Only Trace Model}
Each run produces a sequence of JSONL events. The append-only property supports streaming ingestion, partial recovery from crashes, and deterministic replay without requiring in-memory aggregation.

\subsection{Schema Versioning}
Every trace event includes a required top-level \texttt{schema\_version} field. Breaking changes require a new schema version, preventing silent incompatibility and enabling replay integrity checks.

\subsection{Event Contracts}
Events include required fields \texttt{schema\_version}, \texttt{ts}, \texttt{event}, \texttt{role}, and \texttt{payload}. Event types are explicitly enumerated and include at minimum \texttt{kernel.start}, \texttt{role.planner}, \texttt{role.critic}, \texttt{kernel.revise}, \texttt{kernel.decision}, \texttt{role.executor}, and \texttt{kernel.end}. Each event type defines a stable payload contract within a schema version.

\subsection{Strict Validation}
Validation enforces schema stability and protocol completeness. Non-strict mode permits partial runs; strict mode requires \texttt{kernel.end}, forbids unknown event types, and enforces schema version consistency.

\subsection{Replay and Aggregation}
Replay reconstructs outcomes, revision counts, veto frequencies, override rates, and reason-code distributions without re-executing the original task. Streaming replay emits one summarized JSON object per run for scalable analytics.

\subsection{Summary}
Schema versioning, strict validation, and replay tooling render ROK runs auditable and analyzable over time.

